\begin{theorem}[无外加寄存器、置换自然性约束下的二值跳跃结构唯一性]\label{thm:bdry-binary-jump-orientation-functor-uniqueness}
令 \(k\ge 2\)，并令 \(\mathbf{FinSet}_k^{\cong}\) 表示由所有 \(k\) 元有限集合与其双射组成的群胚（任意两对象同构，且 \(\Aut([k])\cong \mathfrak S_k\)）。
令 \(\mathbf{Tors}(\ZZ_2)\) 表示 \(\ZZ_2\)-torsor（两点集合上的自由传递 \(\ZZ_2\) 作用）与其同构构成的群胚。
称一个\emph{二值跳跃结构}为一函子
\[
\mathcal{J}:\mathbf{FinSet}_k^{\cong}\longrightarrow \mathbf{Tors}(\ZZ_2).
\]
则当 \(k\ge 3\) 时，除平凡函子外，在自然同构意义下存在且仅存在一个非平凡 \(\mathcal{J}\)；
其可由\emph{取向函子}给出：对任意 \(k\) 元集合 \(S\)，定义
\[
\mathrm{Or}(S):=\mathrm{Bij}([k],S)\big/\mathfrak A_k,
\]
其中 \(\mathfrak A_k\subset \mathfrak S_k\) 为交错群，并以 \(\mathfrak S_k\) 在 \(\mathrm{Bij}([k],S)\) 上的预合成作用诱导 \(\mathrm{Or}(S)\) 的 \(\ZZ_2\)-torsor 结构；其作用恰经由符号同态 \(\mathrm{sgn}:\mathfrak S_k\to\ZZ_2\) 因子化。
\leanverified{paper\_bdry\_binary\_jump\_orientation\_functor\_uniqueness}
\end{theorem}

\begin{proof}
取基对象 \(S_0=[k]\)。
任一 \(\ZZ_2\)-torsor 的自同构群自然同构于 \(\ZZ_2\)（由“平移”给出），因此对任意函子 \(\mathcal{J}\)，由 \(\Aut(S_0)=\mathfrak S_k\) 在 \(\mathcal{J}(S_0)\) 上的函子性作用得到群同态
\[
\chi_{\mathcal{J}}:\mathfrak S_k\to \Aut(\mathcal{J}(S_0))\cong \ZZ_2.
\]
反之，给定任意群同态 \(\chi:\mathfrak S_k\to\ZZ_2\)，取一固定两点 \(\ZZ_2\)-torsor \(T\) 并令 \(\mathfrak S_k\) 通过 \(\chi\) 作用于 \(T\)；
再利用 \(\mathbf{FinSet}_k^{\cong}\) 的连通性与函子性，可沿任意双射将该作用在所有对象上传输，从而得到一函子 \(\mathcal{J}\)，且其同构类仅依赖于 \(\chi\)。
由此 \(\mathcal{J}\) 的自然同构类与 \(\mathrm{Hom}(\mathfrak S_k,\ZZ_2)\) 一一对应。
\par
当 \(k\ge 3\) 时 \(\mathfrak S_k^{\mathrm{ab}}\cong \ZZ_2\)，故 \(\mathrm{Hom}(\mathfrak S_k,\ZZ_2)\) 仅含平凡同态与符号同态两者，从而非平凡 \(\mathcal{J}\) 在自然同构意义下唯一。
最后验证 \(\mathrm{Or}\) 对应 \(\mathrm{sgn}\)：\(\mathfrak S_k\) 在 \(\mathrm{Bij}([k],S)\) 上的预合成作用在商集 \(\mathrm{Bij}([k],S)/\mathfrak A_k\) 上恰给出一个两点轨道空间，其作用即由 \(\mathrm{sgn}\) 决定。证毕。
\end{proof}

\begin{corollary}[边界 uplift 的取向 parity：从二层 sheet parity 到三层二值覆盖]\label{cor:bdry-orientation-parity-uplift}
\leanverified{paper\_bdry\_orientation\_parity\_uplift}
沿用 \(\Delta_m^{\mathrm{bdry}}\) 的 boundary uplift 位移集合记号（例如表 \ref{tab:foldbin_boundary_lift_m6_m8} 给出 \(m\in\{6,7,8\}\) 的证书输出），并对每个 boundary 纤维 \(w\) 令 \(\Delta_m^{\mathrm{bdry}}(w)\) 表示其 uplift-delta 作为无序有限集。
在“不引入外加寄存器且对分支重标号保持自然性”的约束下，可定义一个跨层数统一的二值跳跃对象：
\[
\Pi_m(w):=\mathrm{Or}\!\bigl(\Delta_m^{\mathrm{bdry}}(w)\bigr)\in \mathbf{Tors}(\ZZ_2).
\]
则：
\begin{enumerate}
\normalfont
  \item \(\Pi_m(w)\) 对 uplift 分支的任意重标号（置换）天然不变，并随双射函子性变化；
  \item 当 \(|\Delta_m^{\mathrm{bdry}}(w)|=2\) 时（例如 \(m=6\) 的 boundary 二层覆盖），因 \(\mathfrak A_2=\{e\}\) 有规范同构 \(\mathrm{Or}(S)\cong \mathrm{Bij}([2],S)\cong S\)，故 \(\Pi_m(w)\) 退化回二层 sheet parity 的等价表述；
  \item 当 \(|\Delta_m^{\mathrm{bdry}}(w)|=3\) 时（例如 \(m=7,8\) 的三层覆盖），\(\Pi_m(w)\) 给出唯一非平凡、且无需指定“特异层”的二值跳跃对象；其对象类型为取向 torsor，而非“层标签上的 \(\ZZ_2\)-分级”。
\end{enumerate}
\end{corollary}

\begin{proof}
第 (1) 点是 \(\mathrm{Or}\) 的函子性。
第 (2) 点由 \(\mathfrak A_2=\{e\}\) 直接推出。
第 (3) 点由定理 \ref{thm:bdry-binary-jump-orientation-functor-uniqueness} 在 \(k=3\) 的唯一性分类给出：置换自然且非平凡的二值结构在同构意义下只能是取向 torsor。证毕。
\end{proof}

\begin{proposition}[fold-gauge 的符号降维：二值跳跃从中心迁移到交换化]\label{prop:bdry-fold-gauge-sign-abelianization}
对任意有限集合 \(S\)（\(|S|=d\ge 2\)），令
\[
\mathrm{Lab}(S):=\mathrm{Bij}([d],S)
\]
为其标号集合，则 \(\mathrm{Lab}(S)\) 为主 \(\mathfrak S_d\)-torsor。
记 \(\mathrm{sgn}:\mathfrak S_d\to\ZZ_2\) 为符号同态，则存在规范同构
\[
\mathrm{Or}(S)\ \cong\ \mathrm{Lab}(S)\times_{\mathfrak S_d,\mathrm{sgn}}\ZZ_2.
\]
特别地：当 \(d=2\) 时 \(\mathfrak S_2\cong\ZZ_2\) 且其中心即自身，二值跳跃可实现为纤维内的“中心型”层翻转；而当 \(d\ge 3\) 时 \(Z(\mathfrak S_d)=\{e\}\) 但 \(\mathfrak S_d^{\mathrm{ab}}\cong\ZZ_2\)，故二值跳跃只能通过交换化（符号表示）出现，即表现为取向 torsor 而非层标签。
\end{proposition}

\begin{proof}
由定义 \(\mathrm{Or}(S)=\mathrm{Bij}([d],S)/\mathfrak A_d\)。
另一方面，纤维积 \(\mathrm{Lab}(S)\times_{\mathfrak S_d,\mathrm{sgn}}\ZZ_2\) 等价于在 \(\mathrm{Lab}(S)\) 中把相差一个偶置换的标号识别为同类，故与上述商完全一致。
关于中心与交换化的断言为对称群的基本结构：\(Z(\mathfrak S_d)=\{e\}\)（\(d\ge 3\)）且 \(\mathfrak S_d^{\mathrm{ab}}\cong\ZZ_2\)（\(d\ge 2\)）。证毕。
\end{proof}
\leanverified{paper\_bdry\_fold\_gauge\_sign\_abelianization}
\leanverified{paper\_bdry\_fold\_gauge\_sign\_abelianization\_verified}

\begin{proposition}[字符扭曲二值跳跃标记的可实现性判别]\label{prop:bdry-chi-twisted-binary-label-existence}
设有限群 \(G\) 传递作用于有限集合 \(S\)，并取群同态 \(\chi:G\to\ZZ_2\)。
称映射 \(p:S\to\ZZ_2\) 为 \(\chi\)-扭曲二值跳跃标记，若对一切 \(g\in G\)、\(s\in S\) 有
\[
p(g\cdot s)=p(s)+\chi(g)\qquad(\mathrm{mod}\ 2).
\]
则存在 \(\chi\)-扭曲二值跳跃标记当且仅当对某（等价地任意）基点 \(s_0\in S\) 有
\[
\chi|_{G_{s_0}}\equiv 0,
\qquad
G_{s_0}:=\{g\in G:\ g\cdot s_0=s_0\}.
\]
在该条件下，解集在整体翻转 \(p\mapsto p+1\) 下构成一条 \(\ZZ_2\)-torsor。
\end{proposition}
\begin{proof}
必要性：若 \(p\) 满足扭曲条件，则对任意 \(h\in G_{s_0}\) 有
\(
p(s_0)=p(h\cdot s_0)=p(s_0)+\chi(h)
\)，故 \(\chi(h)=0\)。
\par
充分性：若 \(\chi|_{G_{s_0}}\equiv 0\)，定义 \(p(g\cdot s_0):=\chi(g)\)。
若 \(g\cdot s_0=g'\cdot s_0\)，则 \(g^{-1}g'\in G_{s_0}\)，从而 \(\chi(g)=\chi(g')\)，故 \(p\) 良定义。
并且对任意 \(g,g'\in G\) 有
\(
p(gg'\cdot s_0)=\chi(gg')=\chi(g)+\chi(g')=p(g'\cdot s_0)+\chi(g)
\)，代入 \(s=g'\cdot s_0\) 即得扭曲关系。
最后，若 \(p\) 为一解，则 \(p+1\) 仍为解；并且两解之差为常值函数（由传递性与扭曲关系立得），从而解集为 \(\ZZ_2\)-torsor。证毕。
\end{proof}
\leanverified{paper\_bdry\_chi\_twisted\_binary\_label\_existence}

\begin{corollary}[对称群符号特征与二点层标记的刚性]\label{cor:bdry-symmetric-group-sign-twisted-label-d2}
\leanverified{paper\_bdry\_symmetric\_group\_sign\_twisted\_label\_d2}
令 \(\mathfrak S_d\) 自然作用于 \([d]=\{1,\dots,d\}\)，并取 \(\chi=\mathrm{sgn}:\mathfrak S_d\to\ZZ_2\)。
则存在 \(\mathrm{sgn}\)-扭曲二值跳跃标记 \(p:[d]\to\ZZ_2\) 当且仅当 \(d=2\)；
并且当 \(d=2\) 时该标记在整体翻转 \(p\mapsto p+1\) 下唯一。
\end{corollary}
\begin{proof}
由命题 \ref{prop:bdry-chi-twisted-binary-label-existence}，存在性等价于 \(\mathrm{sgn}\) 在稳定子 \((\mathfrak S_d)_{1}\cong \mathfrak S_{d-1}\) 上恒为 \(0\)。
当 \(d\ge 3\) 时 \(\mathfrak S_{d-1}\) 含换位（奇置换），故不可能；当 \(d=2\) 时稳定子平凡，存在性成立。
唯一性由命题 \ref{prop:bdry-chi-twisted-binary-label-existence} 的 torsor 性直接给出。证毕。
\end{proof}

\begin{proposition}[二值分层在置换对称下的轨道分类与精确账本量]\label{prop:bdry-binary-layering-orbit-classification}
令 \(F\) 为 \(d\) 点集合。
定义二值分层的等价类空间
\[
\mathcal P_d(F):=\{p:F\to\ZZ_2\}\big/\bigl(p\sim p+1\bigr),
\]
并令 \(\mathfrak S(F)\cong \mathfrak S_d\) 以置换作用于 \(\mathcal P_d(F)\)。
则轨道空间由单一整数参数完全刻画：
\[
\mathcal P_d(F)\big/\mathfrak S(F)\ \cong\ \{0,1,\dots,\lfloor d/2\rfloor\},
\qquad
k:=\min\bigl(|p^{-1}(0)|,\ |p^{-1}(1)|\bigr).
\]
并且对应轨道大小为
\[
|\mathcal O_k|=
\begin{cases}
\binom{d}{k},& k<d/2,\\[4pt]
\binom{d}{d/2}/2,& d\ \text{为偶且}\ k=d/2.
\end{cases}
\]
\leanverified{paper\_bdry\_binary\_layering\_orbit\_classification}
\end{proposition}
\begin{proof}
对任意 \(p\)，其值域逆像大小为 \((a,d-a)\)；在整体翻转下等价于 \((\min(a,d-a),\max(a,d-a))\)，故轨道不变量可取为 \(k=\min(a,d-a)\)。
反之给定 \(k\le d/2\)，取任意 \(k\)-子集 \(A\subset F\) 并令 \(p=\mathbf 1_A\) 得到代表元。
当 \(k<d/2\) 时，轨道对应于 \(k\)-子集集合，\(\mathfrak S(F)\) 在其上传递，故轨道大小为 \(\binom{d}{k}\)；
当 \(k=d/2\)（\(d\) 偶）时，\(A\) 与补集在 \(\mathcal P_d(F)\) 中同类，故需再除以 \(2\)。证毕。
\end{proof}

\begin{proposition}[取向 torsor 的行列式实现与跳跃类]\label{prop:bdry-orientation-torsor-determinant-realization-jumpclass}
\leanverified{paper\_bdry\_orientation\_torsor\_determinant\_realization\_jumpclass}
设 \(S\) 为有限集合，\(|S|=d\ge 1\)，并令
\[
V(S):=\RR^{S}
\]
为以 \(S\) 为基的实向量空间。
记 \(\mathrm{Or}_{\det}(V(S))\) 为 \(V(S)\) 的取向集合（两点集合，视为 \(\ZZ_2\)-torsor）。
则存在关于 \(S\) 双射自然的规范同构（\(\ZZ_2\)-torsor 意义下）
\[
\Phi_S:\ \mathrm{Or}(S)\ \xrightarrow{\ \sim\ }\ \mathrm{Or}_{\det}(V(S)).
\]
\par
进一步，设 \(Y\) 连通，\(\pi:\widetilde Y\to Y\) 为常度数 \(d\) 的局部有限层覆叠，并令 \(S_y=\pi^{-1}(y)\)。
则 \(y\mapsto \mathrm{Or}(S_y)\) 组装成 \(Y\) 上的 \(\ZZ_2\)-局部系统（主 \(\ZZ_2\)-丛），从而对应一个上同调类
\[
j(\pi)\in H^1(Y;\ZZ_2).
\]
将 \(\pi_\ast\RR\) 视为由纤维置换作用诱导的秩 \(d\) 实局部系统，则有恒等式
\[
j(\pi)=w_1\bigl(\det(\pi_\ast\RR)\bigr)\in H^1(Y;\ZZ_2),
\]
其中 \(\det(\pi_\ast\RR)\) 为行列式线局部系统，\(w_1\) 为第一 Stiefel--Whitney 类。
\end{proposition}
\begin{proof}
取一标号 \(b\in\mathrm{Bij}([d],S)\)，则有序列 \((b(1),\dots,b(d))\) 给出 \(V(S)\) 的一组有序基，从而给出取向。
若以偶置换右乘 \(b\)，基变换的行列式为 \(+1\)，故取向不变。
因此取向仅依赖于 \(b\) 在 \(\mathfrak A_d\) 右陪集中的像，得到良定义映射 \(\Phi_S\)；
其为双射且与任意双射 \(f:S\to S'\) 交换，故为自然同构。
\par
对覆叠 \(\pi\)，沿环路的解析延拓给出纤维的单值化置换 \(\rho(\gamma)\in\mathfrak S_d\)。
由定义 \(\mathrm{Or}(S_y)=\mathrm{Bij}([d],S_y)/\mathfrak A_d\)，其单值化恰通过符号同态 \(\mathrm{sgn}(\rho(\gamma))\in\ZZ_2\)，从而得到 \(j(\pi)\)。
另一方面，局部系统 \(\pi_\ast\RR\) 的单值化为相应置换矩阵表示，其行列式即 \(\mathrm{sgn}(\rho(\gamma))\)，故 \(\det(\pi_\ast\RR)\) 的 \(w_1\) 与 \(j(\pi)\) 表示同一 \(\pi_1(Y)\to\ZZ_2\) 同态，从而在 \(H^1(Y;\ZZ_2)\) 中相等。证毕。
\end{proof}

\begin{corollary}[取向局部系统的同调分类与 \texorpdfstring{$\mathfrak A_d$}{Ad}-约化判别]\label{cor:bdry-orientation-local-system-w1-alternating-reduction}
\leanverified{paper\_bdry\_orientation\_local\_system\_w1\_alternating\_reduction\_structured}
沿用命题 \ref{prop:bdry-orientation-torsor-determinant-realization-jumpclass} 的记号，并设 \(\rho:\pi_1(Y,y_0)\to \mathfrak S_d\) 为覆叠 \(\pi\) 的 monodromy 表示。
则取向局部系统 \(y\mapsto \mathrm{Or}(S_y)\) 的同构类由同态 \(\mathrm{sgn}\circ\rho:\pi_1(Y,y_0)\to\ZZ_2\) 唯一决定，并对应于类 \(j(\pi)\in H^1(Y;\ZZ_2)\)。
特别地，以下命题等价：
\begin{enumerate}
  \item \(j(\pi)=0\in H^1(Y;\ZZ_2)\)；
  \item \(y\mapsto \mathrm{Or}(S_y)\) 存在全局截面（等价地：该 \(\ZZ_2\)-torsor 局部系统为平凡丛）；
  \item \(\rho(\pi_1(Y,y_0))\subseteq \mathfrak A_d\)，即覆叠的结构群可从 \(\mathfrak S_d\) 约化到交错群。
\end{enumerate}
\end{corollary}
\begin{proof}
命题 \ref{prop:bdry-orientation-torsor-determinant-realization-jumpclass} 的证明已给出：\(\mathrm{Or}(S_y)\) 的单值化恰经由 \(\mathrm{sgn}\circ\rho\)，从而其同调类即为 \(j(\pi)\)。
因此 \(j(\pi)=0\) 当且仅当 \(\mathrm{sgn}\circ\rho\) 平凡，当且仅当 \(\rho(\pi_1)\) 落在 \(\ker(\mathrm{sgn})=\mathfrak A_d\)。
另一方面，连通底空间上的主 \(\ZZ_2\)-丛平凡当且仅当存在全局截面，故 (1)\(\Leftrightarrow\)(2)。
证毕。
\end{proof}

\begin{proposition}[纤维积的跳跃类乘法律与偶度稳定化湮灭]\label{prop:bdry-orientation-jumpclass-fiberproduct-multiplicativity}
\leanverified{paper\_bdry\_orientation\_jumpclass\_fiberproduct\_multiplicativity}
设 \(Y\) 连通，\(\pi_i:\widetilde Y_i\to Y\)（\(i=1,2\)）为常度数有限覆叠，\(\deg(\pi_i)=d_i\)。
设纤维积 \(\widetilde Y:=\widetilde Y_1\times_Y\widetilde Y_2\) 在 \(Y\) 上仍为覆叠，并记诱导覆叠为 \(\pi:\widetilde Y\to Y\)，则
\[
j(\pi)\equiv (d_2\bmod 2)\,j(\pi_1)+(d_1\bmod 2)\,j(\pi_2)\qquad\text{于 }H^1(Y;\ZZ_2).
\]
特别地若 \(d_1\) 为偶数，则 \(j(\pi)=(d_2\bmod 2)\,j(\pi_1)\)；
若 \(d_1,d_2\) 皆为偶数，则 \(j(\pi)=0\)。
\end{proposition}
\begin{proof}
由命题 \ref{prop:bdry-orientation-torsor-determinant-realization-jumpclass}，
\(
j(\pi)=w_1(\det(\pi_\ast\RR))
\)。
对常度数覆叠有自然同构
\[
(\pi_\ast\RR)\ \cong\ (\pi_{1\ast}\RR)\otimes(\pi_{2\ast}\RR),
\]
并且对有限维实局部系统 \(U,W\) 有行列式恒等式
\[
\det(U\otimes W)\ \cong\ \det(U)^{\otimes \dim W}\otimes \det(W)^{\otimes \dim U}.
\]
取 \(U=\pi_{1\ast}\RR\)、\(W=\pi_{2\ast}\RR\) 并用 \(w_1(L^{\otimes n})\equiv (n\bmod 2)\,w_1(L)\)（\(\ZZ_2\) 系数）即得所示公式；偶度情形为其直接推论。证毕。
\end{proof}

\begin{theorem}[多重纤维积的取向障碍奇偶公式]\label{thm:bdry-orientation-iterated-fiberproduct-parity}
\leanverified{paper\_bdry\_orientation\_iterated\_fiberproduct\_parity}
设 \(Y\) 连通，\(\pi_i:\widetilde Y_i\to Y\)（\(1\le i\le r\)）为常度数有限覆叠，次数记为 \(d_i\)，并记其取向局部系统障碍类为
\[
j(\pi_i)\in H^1(Y;\ZZ_2).
\]
设
\[
\Pi:=\widetilde Y_1\times_Y\widetilde Y_2\times_Y\cdots\times_Y\widetilde Y_r\longrightarrow Y
\]
为迭代纤维积诱导的覆叠，则有闭式
\[
j(\Pi)=\sum_{i=1}^r\Bigl(\prod_{\ell\neq i} d_\ell \bmod 2\Bigr)\,j(\pi_i)
\qquad\text{于 }H^1(Y;\ZZ_2).
\]
特别地：
\begin{enumerate}
\normalfont
  \item 若至少有两个 \(d_i\) 为偶数，则 \(j(\Pi)=0\)；
  \item 若全部 \(d_i\) 为奇数，则 \(j(\Pi)=\sum_{i=1}^r j(\pi_i)\)；
  \item 若恰有一个 \(d_k\) 为偶数而其余全为奇数，则 \(j(\Pi)=j(\pi_k)\)。
\end{enumerate}
\end{theorem}
\begin{proof}
对 \(r\) 作归纳。\(r=2\) 即命题 \ref{prop:bdry-orientation-jumpclass-fiberproduct-multiplicativity}。
设结论对 \(r-1\) 成立，并记
\[
\Pi_{r-1}:=\widetilde Y_1\times_Y\cdots\times_Y\widetilde Y_{r-1}\to Y,
\qquad
D_{r-1}:=\prod_{i=1}^{r-1}d_i.
\]
由归纳假设，
\[
j(\Pi_{r-1})
=
\sum_{i=1}^{r-1}\Bigl(\prod_{\substack{\ell\neq i\\ \ell\le r-1}} d_\ell \bmod 2\Bigr)\,j(\pi_i).
\]
再将
\(
\Pi_r=\Pi_{r-1}\times_Y\widetilde Y_r
\)
代入命题 \ref{prop:bdry-orientation-jumpclass-fiberproduct-multiplicativity}，得到
\[
j(\Pi_r)
=(d_r\bmod 2)\,j(\Pi_{r-1})+(D_{r-1}\bmod 2)\,j(\pi_r).
\]
把归纳表达式代回，并将 \((d_r\bmod 2)\) 吸收入各项系数，即得
\[
j(\Pi_r)
=
\sum_{i=1}^r\Bigl(\prod_{\ell\neq i} d_\ell \bmod 2\Bigr)\,j(\pi_i).
\]
三条特别情形只是对上述闭式逐项作奇偶判定。证毕。
\end{proof}

\begin{proposition}[笛卡尔积的取向 torsor 分解与指数律]\label{prop:bdry-orientation-cartesian-product-exponent-law}
\leanverified{paper\_bdry\_orientation\_cartesian\_product\_exponent\_law}
对任意有限集合 \(S,T\)，存在关于 \((S,T)\) 双射自然的规范同构
\[
\mathrm{Or}(S\times T)\ \cong\ \mathrm{Or}(S)^{\otimes |T|}\ \otimes\ \mathrm{Or}(T)^{\otimes |S|}.
\]
其中 \(P^{\otimes k}\) 表示在 \(\mathbf{Tors}(\ZZ_2)\) 的张量结构下的 \(k\) 次张量幂（仅依赖 \(k\bmod 2\)）。
\end{proposition}
\begin{proof}
由 \(V(S\times T)\cong V(S)\otimes V(T)\)（基向量按 \((s,t)\) 分解）以及行列式恒等式
\(
\det(U\otimes W)\cong \det(U)^{\otimes \dim W}\otimes \det(W)^{\otimes \dim U}
\)，取 \(U=V(S)\)、\(W=V(T)\) 得
\[
\mathrm{Or}_{\det}(V(S\times T))
\cong
\mathrm{Or}_{\det}(V(S))^{\otimes |T|}\otimes \mathrm{Or}_{\det}(V(T))^{\otimes |S|}.
\]
再用命题 \ref{prop:bdry-orientation-torsor-determinant-realization-jumpclass} 的自然同构 \(\Phi\) 把 \(\mathrm{Or}\) 与 \(\mathrm{Or}_{\det}\) 识别，即得所示同构。证毕。
\end{proof}

\leanverified{paper\_bdry\_orientation\_cartesian\_power\_parity\_collapse}
\begin{theorem}[笛卡尔幂的取向 torsor 幂律与奇偶塌缩]\label{thm:bdry-orientation-cartesian-power-parity-collapse}
\leanverified{paper\_bdry\_orientation\_cartesian\_power\_parity\_collapse}
设 \(S\) 为有限集合，\(|S|=d\ge 1\)，记其取向 torsor 为 \(\mathrm{Or}(S)\)。
则对任意整数 \(r\ge 1\)，存在关于 \(S\) 自然的规范同构
\[
\mathrm{Or}(S^{\times r})\ \cong\ \mathrm{Or}(S)^{\otimes\, r d^{r-1}}.
\]
因此有完整奇偶分类
\[
\mathrm{Or}(S^{\times r})\cong
\begin{cases}
\mathbf 1, & d\ \text{为偶且}\ r\ge 2,\\[2pt]
\mathbf 1, & d\ \text{为奇且}\ r\ \text{为偶},\\[2pt]
\mathrm{Or}(S), & d\ \text{为奇且}\ r\ \text{为奇},
\end{cases}
\]
其中 \(\mathbf 1\) 表示 \(\mathbf{Tors}(\ZZ_2)\) 的单位对象。
\end{theorem}
\begin{proof}
记 \(a_r\) 为使得
\[
\mathrm{Or}(S^{\times r})\cong \mathrm{Or}(S)^{\otimes a_r}
\]
成立的整数。
当 \(r=1\) 时有 \(a_1=1\)。
由命题 \ref{prop:bdry-orientation-cartesian-product-exponent-law}，
\[
\mathrm{Or}(S^{\times(r+1)})
\cong
\mathrm{Or}(S^{\times r})^{\otimes d}\otimes \mathrm{Or}(S)^{\otimes d^r}
\cong
\mathrm{Or}(S)^{\otimes (d a_r+d^r)}.
\]
故 \(a_r\) 满足递推
\[
a_{r+1}=d a_r+d^r,\qquad a_1=1.
\]
直接验证其闭式解为 \(a_r=r d^{r-1}\)，从而得到所示同构。
\par
再看奇偶性。torsor 张量幂只依赖指数模 \(2\)。
若 \(d\) 为偶，则当 \(r\ge 2\) 时 \(r d^{r-1}\) 为偶数，故 \(\mathrm{Or}(S^{\times r})\cong \mathbf 1\)。
若 \(d\) 为奇，则 \(d^{r-1}\equiv 1\pmod 2\)，故
\(
r d^{r-1}\equiv r\pmod 2
\)，
从而得到三分情形。证毕。
\end{proof}

\begin{proposition}[取向函子的对称单子结构与 Koszul 符号]\label{prop:bdry-orientation-functor-symmetric-monoidal-koszul}
令 \(\mathbf{FinSet}^{\cong}\) 为有限集合与双射构成的群胚，张量结构为不交并 \(\sqcup\)。
令 \(\mathbf{Tors}(\ZZ_2)\) 为 \(\ZZ_2\)-torsor 的 Picard 群胚，张量为 torsor 乘积 \(\otimes\)。
则取向函子 \(\mathrm{Or}:\mathbf{FinSet}^{\cong}\to \mathbf{Tors}(\ZZ_2)\) 可提升为对称单子函子：对任意有限集 \(S,T\) 存在自然同构
\[
\mu_{S,T}:\ \mathrm{Or}(S\sqcup T)\xrightarrow{\ \sim\ }\mathrm{Or}(S)\otimes\mathrm{Or}(T),
\]
并且对标准交换双射 \(\beta_{S,T}:S\sqcup T\to T\sqcup S\) 有
\[
\mu_{T,S}\circ \mathrm{Or}(\beta_{S,T})\circ \mu_{S,T}^{-1}
=(-1)^{|S||T|}\cdot\mathrm{id}_{\mathrm{Or}(S)\otimes \mathrm{Or}(T)},
\]
其中 \((-1)\) 表示 \(\mathbf{Tors}(\ZZ_2)\) 中非平凡自同构。
\leanverified{paper\_bdry\_orientation\_functor\_symmetric\_monoidal\_koszul}
\end{proposition}
\begin{proof}
取 \(d=|S|\)、\(e=|T|\)。
任取标号 \(b_S:[d]\xrightarrow{\sim}S\)、\(b_T:[e]\xrightarrow{\sim}T\)，拼接得
\(
b_{S\sqcup T}:[d+e]\xrightarrow{\sim}S\sqcup T
\)
并诱导 \(\mu_{S,T}\)。
交换双射 \(\beta_{S,T}\) 在 \([d+e]\) 上对应把前 \(d\) 个元素整体移到末端的块置换，其奇偶性等于 \((-1)^{de}\)，从而在 \(\mathrm{Or}\) 上诱导 \((-1)^{de}\) 的 torsor 自同构。
自然性与独立于标号选择由定义的商结构直接验证。证毕。
\end{proof}

\begin{proposition}[分块分解的有限集“行列式公式”与奇块可见性]\label{prop:bdry-orientation-block-decomposition-odd-visibility}
\leanverified{paper\_bdry\_orientation\_block\_decomposition\_odd\_visibility}
设有限集合 \(S\) 分解为互不交叠并
\(
S=\bigsqcup_{t\in T}S_t
\)。
令奇块指标子集
\[
T_{\mathrm{odd}}:=\{t\in T:\ |S_t|\equiv 1\ (\mathrm{mod}\ 2)\}\subseteq T.
\]
则在 \(\mathbf{Tors}(\ZZ_2)\) 中存在关于该分块自然的规范同构
\[
\mathrm{Or}(S)\ \cong\ \mathrm{Or}(T_{\mathrm{odd}})\ \otimes\ \bigotimes_{t\in T}\mathrm{Or}(S_t).
\]
换言之，块间重排在取向 torsor 上诱导的二值跳跃仅由奇数大小的块控制：交换两块 \(S_{t_1},S_{t_2}\) 的位置在 \(\mathrm{Or}(S)\) 上引入非平凡翻转当且仅当 \(t_1,t_2\in T_{\mathrm{odd}}\)。
\end{proposition}
\begin{proof}
取 \(T\) 的一标号 \(\theta:[|T|]\xrightarrow{\sim}T\) 与每个块内标号 \(b_t:[|S_t|]\xrightarrow{\sim}S_t\)。
把各块依 \(\theta\) 的顺序拼接（命题 \ref{prop:bdry-orientation-functor-symmetric-monoidal-koszul}）得到 \(S\) 的标号，从而诱导同构
\[
\mathrm{Or}(S)\ \cong\ \bigotimes_{i=1}^{|T|}\mathrm{Or}\!\bigl(S_{\theta(i)}\bigr).
\]
更换 \(\theta\) 的代表仅需考察交换相邻两块 \(S_{t_1},S_{t_2}\)：
该交换在 \([|S|]\) 上对应把一段长度 \(|S_{t_1}|\) 的块穿过一段长度 \(|S_{t_2}|\) 的块，故其符号为 \((-1)^{|S_{t_1}||S_{t_2}|}\)。
因此仅当两块皆为奇大小时该交换在 \(\mathrm{Or}(S)\) 上诱导非平凡翻转。
\par
另一方面，\(\mathrm{Or}(T_{\mathrm{odd}})\) 的对称约束在交换两单点块时恰给出 \((-1)^{1\cdot 1}=-1\)；
而若至少一块为偶，则该交换不作用于 \(T_{\mathrm{odd}}\)，从而在 \(\mathrm{Or}(T_{\mathrm{odd}})\) 上平凡。
由此可见：将 \(\mathrm{Or}(T_{\mathrm{odd}})\) 作为额外因子后，上述拼接同构对 \(\theta\) 的改变严格不变，因而得到与标号无关的规范同构
\(
\mathrm{Or}(S)\cong \mathrm{Or}(T_{\mathrm{odd}})\otimes\bigotimes_{t\in T}\mathrm{Or}(S_t)
\)。
证毕。
\end{proof}

\begin{theorem}[奇块可见性定律与 wreath 跳跃字符的完全公式]\label{thm:bdry-orientation-wreath-character-odd-visibility}
\leanverified{paper\_bdry\_orientation\_wreath\_character\_odd\_visibility}
设有限集合按分块分解为
\[
S=\bigsqcup_{t\in T}S_t,
\qquad
T_{\mathrm{odd}}:=\{t\in T:\ |S_t|\equiv 1\pmod 2\}.
\]
沿用命题 \ref{prop:bdry-orientation-block-decomposition-odd-visibility} 的规范同构
\[
\mathrm{Or}(S)\cong \mathrm{Or}(T_{\mathrm{odd}})\otimes \bigotimes_{t\in T}\mathrm{Or}(S_t).
\]
设 \(g\in \mathrm{Sym}(S)\) 保持该分块，即 \(g\) 诱导一个分块置换 \(\tau\in \mathrm{Sym}(T)\)，并对每个 \(t\in T\) 诱导双射
\[
g_t:\ S_t\xrightarrow{\sim} S_{\tau(t)}.
\]
记 \(\varepsilon(g_t)\in\{\pm 1\}\) 为 \(g_t\) 在取向 torsor 上诱导的符号，即在任意标号识别下 \(\mathrm{Or}(g_t)\) 对应于符号 \(\varepsilon(g_t)\)。
则 \(g\) 在 \(\mathrm{Or}(S)\) 上诱导的字符满足
\[
\varepsilon_S(g)
=
\mathrm{sgn}\!\bigl(\tau|_{T_{\mathrm{odd}}}\bigr)\cdot
\prod_{t\in T}\varepsilon(g_t).
\]

特别地，若所有块大小都相同且等于 \(d\)，并固定一个 \(d\) 元模型块 \(S_d\)，从而把 \(S\) 识别为 \(T\times S_d\)，则对 wreath 乘积元素
\[
(\sigma_t)_{t\in T}\rtimes \tau\in \mathrm{Sym}(S_d)\wr \mathrm{Sym}(T)
\]
有
\[
\varepsilon_S\bigl((\sigma_t)\rtimes \tau\bigr)
=
\mathrm{sgn}(\tau)^d\cdot \prod_{t\in T}\mathrm{sgn}(\sigma_t).
\]
因而当 \(d\) 为偶数时，顶层置换 \(\tau\) 在取向层面完全不可见；当 \(d\) 为奇数时，顶层置换的奇偶性完整保留。
\end{theorem}
\begin{proof}
由命题 \ref{prop:bdry-orientation-block-decomposition-odd-visibility}，
\(\mathrm{Or}(S)\) 分解为奇块标签 torsor 与各块内部取向 torsor 的张量积。
分块置换 \(\tau\) 对 \(\mathrm{Or}(T_{\mathrm{odd}})\) 的作用正是其在奇块指标集上的符号表示，因此贡献
\(
\mathrm{sgn}(\tau|_{T_{\mathrm{odd}}})
\)。
另一方面，每个 \(g_t:S_t\to S_{\tau(t)}\) 在取向 torsor 上诱导 \(\varepsilon(g_t)\)；将这些因子张量相乘即得总字符
\[
\varepsilon_S(g)
=
\mathrm{sgn}\!\bigl(\tau|_{T_{\mathrm{odd}}}\bigr)\cdot
\prod_{t\in T}\varepsilon(g_t).
\]
\par
当所有块大小都等于 \(d\) 时，若 \(d\) 为偶，则 \(T_{\mathrm{odd}}=\varnothing\)，顶层置换贡献为 \(1\)；
若 \(d\) 为奇，则 \(T_{\mathrm{odd}}=T\)，顶层置换贡献为 \(\mathrm{sgn}(\tau)\)。
这恰可统一写成 \(\mathrm{sgn}(\tau)^d\)。
而此时每个 \(g_t\) 由模型块上的 \(\sigma_t\in \mathrm{Sym}(S_d)\) 表示，其取向字符即 \(\mathrm{sgn}(\sigma_t)\)，从而得到 wreath 情形的公式。证毕。
\end{proof}

\begin{proposition}[\texorpdfstring{$t$}{t}-子集构形覆叠的取向传递系数]\label{prop:bdry-orientation-tsubset-transfer-binomial}
设 \(S\) 为有限集合，\(|S|=k\ge 2\)，并记
\[
\binom{S}{t}:=\{A\subset S:\ |A|=t\}\qquad(1\le t\le k-1).
\]
则存在关于 \((S,t)\) 自然的规范同构
\[
\mathrm{Or}\!\left(\binom{S}{t}\right)\ \cong\ \mathrm{Or}(S)^{\otimes\,\binom{k-2}{t-1}}.
\]
等价地，对任意 \(\sigma\in \mathfrak S_k\) 的自然作用有
\[
\mathrm{sgn}\bigl(\sigma\curvearrowright \tbinom{S}{t}\bigr)=\mathrm{sgn}(\sigma)^{\binom{k-2}{t-1}}.
\]
\end{proposition}
\begin{proof}
由于 \(\mathfrak S_k^{\mathrm{ab}}\cong \ZZ_2\)，对称群在任何有限集上的符号特征仅可能为平凡或 \(\mathrm{sgn}\) 的幂，因此只需计算一个换位的作用。
取换位 \(\tau\in\mathfrak S_k\)。
在 \(t\)-子集上，\(\tau\) 的不动点是同时含被交换两元或同时不含两元的子集；其余子集按“恰含其一”成对互换。
而“含第一元不含第二元”的子集数为 \(\binom{k-2}{t-1}\)，故 \(\tau\) 在 \(\binom{S}{t}\) 上分解为 \(\binom{k-2}{t-1}\) 个互不相交的换位，从而
\(
\mathrm{sgn}(\tau\curvearrowright \binom{S}{t})=(-1)^{\binom{k-2}{t-1}}
\)。
于是整体符号特征为 \(\mathrm{sgn}^{\binom{k-2}{t-1}}\)，与定理 \ref{thm:bdry-binary-jump-orientation-functor-uniqueness} 的唯一性分类相容，得到所示同构。证毕。
\end{proof}

\leanverified{paper\_bdry\_orientation\_tsubset\_transfer\_binomial}

\leanverified{paper\_bdry\_orientation\_tsubset\_lucas\_parity\_seeds}
\leanverified{paper\_bdry\_orientation\_tsubset\_lucas\_parity\_package}
\leanverified{paper\_bdry\_orientation\_tsubset\_lucas\_parity}
\begin{corollary}[Lucas 奇偶判据与 \texorpdfstring{$t$}{t}-子集操作的跳跃保留/湮灭]\label{cor:bdry-orientation-tsubset-lucas-parity}
沿用命题 \ref{prop:bdry-orientation-tsubset-transfer-binomial} 的记号，\(\binom{k-2}{t-1}\) 为奇数当且仅当 \(t-1\) 的二进制展开逐位不超过 \(k-2\)（Lucas 定理的模 \(2\) 判据）。
因此对任意常度数覆叠 \(\pi:\widetilde Y\to Y\)，若令 \(\pi^{(t)}\) 为诱导的 \(t\)-子集构形覆叠，则
\[
j(\pi^{(t)})\equiv \binom{k-2}{t-1}\,j(\pi)\qquad\text{于 }H^1(Y;\ZZ_2),
\]
从而 \(\binom{k-2}{t-1}\) 偶时诱导跳跃类必湮灭，奇时与原跳跃类同类。
\end{corollary}
\begin{proof}
第一句为 Lucas 定理在模 \(2\) 的标准推论。
第二句由命题 \ref{prop:bdry-orientation-tsubset-transfer-binomial} 与命题 \ref{prop:bdry-orientation-torsor-determinant-realization-jumpclass} 立即得到：诱导覆叠的单值化置换作用把符号特征提升为 \(\mathrm{sgn}^{\binom{k-2}{t-1}}\)，从而在 \(H^1(Y;\ZZ_2)\) 中乘以同一奇偶系数。证毕。
\end{proof}

\begin{proposition}[显式 \texorpdfstring{$t$}{t}-层选择的最小同变寄存器度数]\label{prop:bdry-equvariant-register-min-degree-choose}
固定 \(k\ge 3\) 与 \(1\le t\le k-1\)，并记 \(C_t:=\binom{[k]}{t}\) 为传递的 \(\mathfrak S_k\)-集合。
设 \(R\) 为一有限传递 \(\mathfrak S_k\)-集合，并设存在 \(\mathfrak S_k\)-同变映射 \(f:R\to C_t\)。
则必有
\[
|R|\ \ge\ |C_t|\ =\ \binom{k}{t},
\]
且该下界可达：取 \(R=C_t\)、\(f=\mathrm{id}\) 即取等。
因此任何在置换自然性约束下实现显式 \(t|(k-t)\) 分级的机制，都至少需要 \(\log_2\binom{k}{t}\) 比特的外加寄存器容量。
\end{proposition}
\begin{proof}
由传递性可写 \(R\cong \mathfrak S_k/H\)、\(C_t\cong \mathfrak S_k/H_t\)，其中 \(H\) 为某一稳定子，且 \(H_t\cong \mathfrak S_t\times \mathfrak S_{k-t}\)。
存在 \(\mathfrak S_k\)-同变映射 \(\mathfrak S_k/H\to \mathfrak S_k/H_t\) 当且仅当 \(H\) 共轭包含于 \(H_t\)；
因此 \([\,\mathfrak S_k:H\,]\ge [\,\mathfrak S_k:H_t\,]=\binom{k}{t}\)，即得 \(|R|\ge \binom{k}{t}\)。
取 \(H=H_t\) 得可达性；比特下界为取对数的直接翻译。证毕。
\end{proof}

\leanverified{paper\_bdry\_equvariant\_register\_min\_degree\_choose}

\endinput
